\section{Convergence Analysis}

\label{sec:convergence}

\subsection{Aggregation Convergence}

We prove that the aggregated prior converges to the honest-majority consensus
as the number of participants grows.

\begin{theorem}[Aggregation Convergence]
\label{thm:convergence}
Let $\Delta^*$ be the true optimal prior for the task distribution. Under
Assumption~\ref{ass:threat}, with $n$ total participants ($n - f$ honest), the
aggregated prior satisfies:
\begin{equation}
    \mathbb{E}\left[\|\tilde{\Delta} - \Delta^*\|_1\right]
    \leq \frac{C}{\sqrt{n - f}} + O\left(\frac{f}{n}\right),
\end{equation}
where $C$ depends on the variance of honest submissions.
\end{theorem}

\begin{proof}
Decompose the error into honest estimation error and Byzantine corruption:
\begin{equation}
    \|\tilde{\Delta} - \Delta^*\|_1
    \leq \underbrace{\|\tilde{\Delta} - \Delta_H^*\|_1}_{\text{Byzantine deviation}}
    + \underbrace{\|\Delta_H^* - \Delta^*\|_1}_{\text{honest estimation error}},
\end{equation}
where $\Delta_H^*$ is the weighted median of honest submissions only.

\textit{Byzantine deviation:} By Theorem~\ref{thm:byz-robust}, the weighted
median with Byzantine nodes equals the weighted median of a subset of honest
values. The deviation from the full honest median is bounded by the gap between
adjacent honest submissions, which is $O(1/\sqrt{n-f})$ by concentration.

\textit{Honest estimation error:} The honest median converges to $\Delta^*$ at
rate $O(1/\sqrt{n-f})$ by standard median convergence for bounded random
variables~\cite{minsker2015geometric}.

The $O(f/n)$ term accounts for the influence of the quality gate: some honest
priors near the quality boundary may be incorrectly filtered when Byzantine
nodes manipulate quality votes.
\end{proof}

\subsection{Prior Quality Dynamics}

The quality of the aggregated prior bank evolves over time as new submissions
arrive and old ones are pruned.

\begin{proposition}[Quality Monotonicity]
Under the condition that the expected quality of new submissions exceeds the
current aggregate quality:
\begin{equation}
    \mathbb{E}[q_{\text{new}}] > \tilde{q}_{\text{current}},
\end{equation}
the aggregate quality $\tilde{q}$ is non-decreasing in expectation over
successive aggregation rounds.
\end{proposition}

\subsection{Network Effect}

As more agents join the network and contribute priors, each individual agent
benefits from the collective knowledge:

\begin{corollary}[Network Scaling]
The expected performance improvement for each agent scales as:
\begin{equation}
    \Delta_{\text{perf}}(n) = \Theta\left(\sqrt{\frac{n-f}{n_0}}\right) \cdot \Delta_{\text{perf}}(n_0),
\end{equation}
where $n_0$ is the baseline number of agents and $\Delta_{\text{perf}}(n_0)$ is
the improvement with $n_0$ agents. This sub-linear scaling reflects diminishing
returns from additional priors as the prior bank becomes saturated.
\end{corollary}

% ============================================================================
% 9. INCENTIVE MECHANISMS
% ============================================================================
